% GENERATED FILE. Edit canonical lessons, then run npm run generate:books.
% canonical-source-hash: fnv1a64:575497dbce57eb00
% canonical-lessons: ZH-C528-pengqiao, ZH-C528-budebu, ZH-C528-chouqian, ZH-C528-dadu, ZH-C528-ouran

\chapter{By Chance}
\label{ch:zh-by-chance}
\hlchaptermodality{528}

\begin{chapteropening}
\textbf{By the end of this chapter:} \emph{I can say something happens by chance, that I have no choice, and talk about drawing lots and making bets.}

Complete the last lesson of chapter 528: I can say something happens by chance, that I have no choice, and talk about drawing lots and making bets.
\end{chapteropening}

\section[pèngqiǎo]{pèngqiǎo — as luck would have it; coincidentally}

\label{lesson:ZH-C528-pengqiao}



\begin{quote}
Before the new one: say the Mandarin for a favourable deal, then the Mandarin for a checked pattern.
\end{quote}

\subsection*{You'll want to know: pèngqiǎo}
\emph{pèngqiǎo} — \textquotedblleft{}as luck would have it; coincidentally\textquotedblright{}.

\begin{quote}
\emph{yàoshi pèngqiǎo lùguò jiù lái wǒ jiā zuòzuò} — If you happen to be passing, come round to my place.
\end{quote}

Its characters come later; for now it is a word for the ear.

\begin{cousinweb}
\emph{pèngqiǎo} is built from \emph{pèng} 'to bump into' and \emph{qiǎo}. No dictionary we use records an older form of the first, so its history is not settled. The second is inherited from Middle Chinese \emph{khaewX} (Guangyun, in Baxter's transcription), glossed 'fine' and 'skilful'; the sense 'by coincidence' is not given there.
\end{cousinweb}

\begin{grammarlens}[title={What you've built}]
One more word for saying what would happen.
\end{grammarlens}

\subsection*{Your turn}
\begin{itemize}
\raggedright
  \item \emph{Say it:} \emph{pèngqiǎo}
  \item \emph{Say it:} \emph{pèngqiǎo}, once more
  \item \emph{Say it:} \emph{yàoshi pèngqiǎo lùguò jiù lái wǒ jiā zuòzuò} — If you happen to be passing, come round to my place.
  \item \emph{Recall:} say the Mandarin for a favourable deal, then the Mandarin for a checked pattern, then say \emph{pèngqiǎo} again
\end{itemize}

\begin{tcolorbox}[breakable,colback=teal!4,colframe=teal!35!black,title={Before you move on}]
What does pèngqiǎo mean? (\textquotedblleft{}As luck would have it; coincidentally\textquotedblright{}.) Say it once more.
\end{tcolorbox}
\section[\texorpdfstring{\zh{不得不} (bùdébù)}{不得不 (bùdébù)}]{\zh{不得不} (bùdébù) — to be obliged to; cannot help but}

\label{lesson:ZH-C528-budebu}



\begin{quote}
Before the new one: say the Mandarin for a checked pattern, then the Mandarin for as luck would have it.
\end{quote}

\subsection*{You'll want to know: \zh{不得不}}
\textbf{\zh{不得不}} — \emph{bùdébù} — \textquotedblleft{}to be obliged to; cannot help but\textquotedblright{}.

\begin{quote}
\textbf{\zh{如果明天下雨我们就不得不改期}} \emph{rúguǒ míngtiān xiàyǔ wǒmen jiù bùdébù gǎiqī} — If it rains tomorrow, we'll have to change the date.
\end{quote}

Tone: the first \zh{不} keeps its falling tone before dé; the second is said bú before a falling tone (bùdébú qù) and bù before the others.

\begin{cousinweb}
\textbf{\zh{不得不}} (\emph{bùdébù}) is built from \zh{不} 'not', \zh{得} 'to get' and \zh{不} again, all inherited from Middle Chinese, \emph{pjut}, \emph{tok} and \emph{pjut} (Guangyun, in Baxter's transcription): may not not do it.
\end{cousinweb}

\begin{grammarlens}[title={What you've built}]
One more word for saying what would happen.
\end{grammarlens}

\subsection*{Your turn}
\begin{itemize}
\raggedright
  \item \emph{Say it:} \emph{bùdébù}
  \item \emph{Say it:} \emph{bùdébù}, once more
  \item \emph{Say it:} \emph{rúguǒ míngtiān xiàyǔ wǒmen jiù bùdébù gǎiqī} — If it rains tomorrow, we'll have to change the date.
  \item \emph{Recall:} say the Mandarin for a checked pattern, then the Mandarin for as luck would have it, then say \emph{bùdébù} again
\end{itemize}

\begin{tcolorbox}[breakable,colback=teal!4,colframe=teal!35!black,title={Before you move on}]
What does \zh{不得不} mean? (\textquotedblleft{}To be obliged to; cannot help but\textquotedblright{}.) Say it once more.
\end{tcolorbox}
\section[chōuqiān]{chōuqiān — to draw lots}

\label{lesson:ZH-C528-chouqian}



\begin{quote}
Before the new one: say the Mandarin for as luck would have it, then the Mandarin for to be obliged to.
\end{quote}

\subsection*{You'll want to know: chōuqiān}
\emph{chōuqiān} — \textquotedblleft{}to draw lots\textquotedblright{}.

\begin{quote}
\emph{yàoshi dàjiā dōu xiǎng qù wǒmen jiù chōuqiān} — If everyone wants to go, let's draw lots.
\end{quote}

Its characters come later; for now it is a word for the ear.

\begin{cousinweb}
\emph{chōuqiān} is built from \emph{chōu} 'to pull out' and \emph{qiān}, both inherited from Middle Chinese, \emph{trhjuw} and \emph{tshjem} (Guangyun, in Baxter's transcription), which glosses the first 'to draw out' and the second, after an older dictionary, 'a test' and 'pointed'; the slip that is drawn is not named there.
\end{cousinweb}

\begin{grammarlens}[title={What you've built}]
One more word for saying what would happen.
\end{grammarlens}

\subsection*{Your turn}
\begin{itemize}
\raggedright
  \item \emph{Say it:} \emph{chōuqiān}
  \item \emph{Say it:} \emph{chōuqiān}, once more
  \item \emph{Say it:} \emph{yàoshi dàjiā dōu xiǎng qù wǒmen jiù chōuqiān} — If everyone wants to go, let's draw lots.
  \item \emph{Recall:} say the Mandarin for as luck would have it, then the Mandarin for to be obliged to, then say \emph{chōuqiān} again
\end{itemize}

\begin{tcolorbox}[breakable,colback=teal!4,colframe=teal!35!black,title={Before you move on}]
What does chōuqiān mean? (\textquotedblleft{}To draw lots\textquotedblright{}.) Say it once more.
\end{tcolorbox}
\section[dǎdǔ]{dǎdǔ — to make a bet}

\label{lesson:ZH-C528-dadu}



\begin{quote}
Before the new one: say the Mandarin for to be obliged to, then the Mandarin for to draw lots.
\end{quote}

\subsection*{You'll want to know: dǎdǔ}
\emph{dǎdǔ} — \textquotedblleft{}to make a bet\textquotedblright{}.

\begin{quote}
\emph{wǒ gēn nǐ dǎdǔ tā míngtiān bú huì lái} — I bet you he won't come tomorrow.
\end{quote}

Its characters come later; for now it is a word for the ear.

Two third tones meet in this word: dǎ is said with a rising tone before dǔ, and the spelling keeps both marks.

\begin{cousinweb}
\emph{dǎdǔ} is built from \zh{打} and \emph{dǔ} 'to gamble', both inherited from Middle Chinese, \emph{taengX} and \emph{tuX} (Guangyun, in Baxter's transcription), which glosses the first 'to strike' and the second as gambling for play. The vowel of dǎ is not the regular outcome of \emph{taengX}.
\end{cousinweb}

\begin{grammarlens}[title={What you've built}]
One more word for saying what would happen.
\end{grammarlens}

\subsection*{Your turn}
\begin{itemize}
\raggedright
  \item \emph{Say it:} \emph{dǎdǔ}
  \item \emph{Say it:} \emph{dǎdǔ}, once more
  \item \emph{Say it:} \emph{wǒ gēn nǐ dǎdǔ tā míngtiān bú huì lái} — I bet you he won't come tomorrow.
  \item \emph{Recall:} say the Mandarin for to be obliged to, then the Mandarin for to draw lots, then say \emph{dǎdǔ} again
\end{itemize}

\begin{tcolorbox}[breakable,colback=teal!4,colframe=teal!35!black,title={Before you move on}]
What does dǎdǔ mean? (\textquotedblleft{}To make a bet\textquotedblright{}.) Say it once more.
\end{tcolorbox}
\section[ǒurán]{ǒurán — accidental; happening by chance}

\label{lesson:ZH-C528-ouran}



\begin{quote}
Before the new one: say the Mandarin for to draw lots, then the Mandarin for to make a bet.
\end{quote}

\subsection*{You'll want to know: ǒurán}
\emph{ǒurán} — \textquotedblleft{}accidental; happening by chance\textquotedblright{}.

\begin{quote}
\emph{rúguǒ zhǐ fāshēng yí cì kěnéng shì ǒurán} — If it only happens once, it may be chance.
\end{quote}

Its characters come later; for now it is a word for the ear.

\begin{cousinweb}
\emph{ǒurán} is built from \emph{ǒu} and \zh{然} 'so', both inherited from Middle Chinese, \emph{nguwH} and \emph{nyen} (Guangyun, in Baxter's transcription). The dictionary glosses \emph{nguwH} 'unexpected', and lists a rising-tone reading, \emph{nguwX}, glossed 'a pair'; the rising tone of ǒu matches that other reading.
\end{cousinweb}

\begin{grammarlens}[title={What you've built}]
One more word for saying what would happen.
\end{grammarlens}

\subsection*{Your turn}
\begin{itemize}
\raggedright
  \item \emph{Say it:} \emph{ǒurán}
  \item \emph{Say it:} \emph{ǒurán}, once more
  \item \emph{Say it:} \emph{rúguǒ zhǐ fāshēng yí cì kěnéng shì ǒurán} — If it only happens once, it may be chance.
  \item \emph{Recall:} say the Mandarin for to draw lots, then the Mandarin for to make a bet, then say \emph{ǒurán} again
\end{itemize}

\begin{tcolorbox}[breakable,colback=teal!4,colframe=teal!35!black,title={Before you move on}]
What does ǒurán mean? (\textquotedblleft{}Accidental; happening by chance\textquotedblright{}.) Say it once more.
\end{tcolorbox}
